\documentclass[../../main]{subfiles}

\begin{document}

\section{Operações sobre Funções}
\label{sec:matematica:funcoes:operacoes}

\subsection{Operações Induzidas}

\begin{defn}[Operação Induzida]
  \label{def:matematica:funcoes:operacoes:induzida}

  Sejam

  \[
    A
  \]

  um conjunto e

  \[
    (B,\star)
  \]

  um conjunto munido de uma operação binária.

  Define-se a operação induzida em

  \[
    B^A
  \]

  por

  \[
    (f\star g)(x)
    =
    f(x)\star g(x)
  \]

  para todo

  \[
    x\in A.
  \]

\end{defn}

\begin{prop}
  \label{prop:matematica:funcoes:operacoes:boa-definicao}

  A operação induzida está bem definida.

\end{prop}

\begin{proof}

  Para cada

  \[
    x\in A,
  \]

  temos

  \[
    f(x),g(x)\in B.
  \]

  Como

  \[
    \star:B\times B\to B,
  \]

  segue que

  \[
    f(x)\star g(x)\in B.
  \]

  Portanto

  \[
    f\star g:A\to B.
  \]

\end{proof}

\subsection{Igualdade Ponto a Ponto}

\begin{prop}[Princípio da Extensionalidade]
  \label{prop:matematica:funcoes:operacoes:extensionalidade}

  Sejam

  \[
    f,g:A\to B.
  \]

  Então

  \[
    f=g
  \]

  se, e somente se,

  \[
    (\forall x\in A)
    \;
    f(x)=g(x).
  \]

\end{prop}

\begin{proof}

  Trata-se da definição de igualdade de funções.

\end{proof}

\subsection{Soma de Funções}

\begin{defn}[Soma]
  \label{def:matematica:funcoes:operacoes:soma}

  Seja

  \[
    (B,+)
  \]

  um conjunto com operação de soma.

  Para

  \[
    f,g\in B^A,
  \]

  define-se

  \[
    f+g
  \]

  por

  \[
    (f+g)(x)
    =
    f(x)+g(x).
  \]

\end{defn}

\begin{ex}

  Se

  \[
    f(x)=x
  \]

  e

  \[
    g(x)=x^2,
  \]

  então

  \[
    (f+g)(x)
    =
    x+x^2.
  \]

\end{ex}

\subsection{Produto de Funções}

\begin{defn}[Produto]
  \label{def:matematica:funcoes:operacoes:produto}

  Seja

  \[
    (B,\cdot)
  \]

  um conjunto com operação multiplicativa.

  Define-se

  \[
    (fg)(x)
    =
    f(x)g(x).
  \]

\end{defn}

\begin{ex}

  Se

  \[
    f(x)=x
  \]

  e

  \[
    g(x)=x^2,
  \]

  então

  \[
    (fg)(x)
    =
    x^3.
  \]

\end{ex}

\subsection{Multiplicação por Escalar}

\begin{defn}[Multiplicação por Escalar]
  \label{def:matematica:funcoes:operacoes:escalar}

  Seja

  \[
    K
  \]

  um corpo.

  Para

  \[
    \lambda\in K
  \]

  e

  \[
    f:A\to K,
  \]

  define-se

  \[
    \lambda f
  \]

  por

  \[
    (\lambda f)(x)
    =
    \lambda\,f(x).
  \]

\end{defn}

\begin{ex}

  Se

  \[
    f(x)=x^2,
  \]

  então

  \[
    (3f)(x)
    =
    3x^2.
  \]

\end{ex}

\subsection{Ordem Ponto a Ponto}

\begin{defn}[Ordem Ponto a Ponto]
  \label{def:matematica:funcoes:operacoes:ordem}

  Seja

  \[
    (B,\le)
  \]

  um conjunto ordenado.

  Para

  \[
    f,g\in B^A,
  \]

  define-se

  \[
    f\le g
  \]

  quando

  \[
    f(x)\le g(x)
  \]

  para todo

  \[
    x\in A.
  \]

\end{defn}

\begin{prop}
  \label{prop:matematica:funcoes:operacoes:ordem-parcial}

  Se

  \[
    (B,\le)
  \]

  é um conjunto parcialmente ordenado, então

  \[
    (B^A,\le)
  \]

  também é parcialmente ordenado.

\end{prop}

\begin{proof}

  As propriedades de reflexividade, antissimetria e transitividade são
  verificadas ponto a ponto.

\end{proof}

\subsection{Herança de Estruturas}

\begin{prop}
  \label{prop:matematica:funcoes:operacoes:grupo}

  Se

  \[
    (B,\star)
  \]

  é um grupo, então

  \[
    (B^A,\star)
  \]

  também é um grupo.

\end{prop}

\begin{proof}

  Todas as propriedades são herdadas ponto a ponto.

\end{proof}

\begin{prop}
  \label{prop:matematica:funcoes:operacoes:anel}

  Se

  \[
    B
  \]

  é um anel, então

  \[
    B^A
  \]

  é um anel com as operações ponto a ponto.

\end{prop}

\begin{proof}

  As propriedades seguem diretamente das propriedades de \(B\).

\end{proof}

\begin{prop}
  \label{prop:matematica:funcoes:operacoes:espaco-vetorial}

  Se

  \[
    K
  \]

  é um corpo, então

  \[
    K^A
  \]

  é um espaço vetorial.

\end{prop}

\begin{proof}

  Os axiomas de espaço vetorial são satisfeitos ponto a ponto.

\end{proof}

\end{document}
